\labsection{8}{BSTs, heaps, and balancing structures}{sec:lab08}

\begin{sourcealignment}{Official Lab Manual --- Lab VIII}
BST operations, min/max heaps, height balancing, and instructor-led B-tree practice.
\end{sourcealignment}

\begin{labproject}{Compare ordering rules across specialized trees}
\textbf{Goal.} Compare BST ordering, heap ordering, and the effect of tree height.

\textbf{Recall first.} Review the \hyperref[ch:trees-ii]{BSTs, Heaps, and Height Balancing} chapter.

\textbf{Complete source files.} Each program below is a separate C++17 file.
Copy its code blocks in printed order, including headers, definitions, helpers, and \texttt{main()}.
The line ranges refer to the supplied file; continued blocks preserve its braces and scope.
Build one program at a time from its \texttt{all\_experiments/Lab08} folder.

\textbf{Program 1: BST operations and height checks.}

\textbf{File:} \texttt{bst.cpp}.

\textbf{Chapter connection:} BST insertion; BST deletion; Height-balancing trees.

Reuse the chapter comparison rule for both insertion and search. Deletion returns the replacement subtree, so callers must assign its result. main covers leaf, one-child, and two-child deletion, then contrasts a chain with a shorter tree.

\begin{adslabfigure}[How this program runs]
\begin{adslabgraphic}
\begin{tikzpicture}[font=\small,execute at begin node={\hyphenpenalty=10000\exhyphenpenalty=10000}]
\node[utpstep,text width=2.50cm,minimum height=1.15cm] (s0) at (0.0,0) {\textbf{1}\\Insert by ordering};
\node[utpstep,text width=2.50cm,minimum height=1.15cm] (s1) at (2.95,0) {\textbf{2}\\Search one branch};
\node[utpstep,text width=2.50cm,minimum height=1.15cm] (s2) at (5.9,0) {\textbf{3}\\Delete / return subtree};
\node[utpstep,text width=2.50cm,minimum height=1.15cm] (s3) at (8.85,0) {\textbf{4}\\Check height / free};
\draw[utparrow] (s0.east)--(s1.west);
\draw[utparrow] (s1.east)--(s2.west);
\draw[utparrow] (s2.east)--(s3.west);
\end{tikzpicture}
\end{adslabgraphic}
\end{adslabfigure}


\begin{utpsteps}\setcounter{enumi}{0}\item \textbf{Headers and initial definitions.}\end{utpsteps}

\texttt{bst.cpp} --- lines 14--30.

\lstinputlisting[firstline=14,lastline=30]{all_experiments/Lab08/bst.cpp}


\tcbbreak

\begin{utpsteps}\setcounter{enumi}{1}\item \textbf{Implement insertBST, contains, deleteBST.}\end{utpsteps}

\texttt{bst.cpp} --- lines 31--65.

% Worked lab input; chapter fonts, colours and TikZ styles retained.
\begin{utpcompactlabfigure}[Three deletion cases applied to the same BST]
\begin{labgraphic}
\begin{tikzpicture}[font=\small]
\begin{scope}[xshift=0cm,yshift=0cm]
\node[utplabel,anchor=south,color=UTPNavy] at (0,.65) {Start};
\node[utpnode,minimum size=7mm] (n0) at (0,0) {50};
\node[utpnode,minimum size=7mm] (n1) at (-1.4,-1) {30};
\node[utpnode,minimum size=7mm] (n2) at (1.4,-1) {70};
\node[utpnodehot,minimum size=7mm] (n3) at (-2.1,-2) {20};
\node[utpnode,minimum size=7mm] (n4) at (-0.7,-2) {40};
\node[utpnode,minimum size=7mm] (n5) at (0.7,-2) {60};
\node[utpnode,minimum size=7mm] (n6) at (2.1,-2) {80};
\draw[utpedge] (n0)--(n1);
\draw[utpedge] (n0)--(n2);
\draw[utpedge] (n1)--(n3);
\draw[utpedge] (n1)--(n4);
\draw[utpedge] (n2)--(n5);
\draw[utpedge] (n2)--(n6);
\end{scope}\begin{scope}[xshift=5.8cm,yshift=0cm]
\node[utplabel,anchor=south,color=UTPNavy] at (0,.65) {1. Delete leaf 20};
\node[utpnode,minimum size=7mm] (n0) at (0,0) {50};
\node[utpnodehot,minimum size=7mm] (n1) at (-1.4,-1) {30};
\node[utpnode,minimum size=7mm] (n2) at (1.4,-1) {70};
\node[utpnode,minimum size=7mm] (n4) at (-0.7,-2) {40};
\node[utpnode,minimum size=7mm] (n5) at (0.7,-2) {60};
\node[utpnode,minimum size=7mm] (n6) at (2.1,-2) {80};
\draw[utpedge] (n0)--(n1);
\draw[utpedge] (n0)--(n2);
\draw[utpedge] (n1)--(n4);
\draw[utpedge] (n2)--(n5);
\draw[utpedge] (n2)--(n6);
\end{scope}\begin{scope}[xshift=0cm,yshift=-3.8cm]
\node[utplabel,anchor=south,color=UTPNavy] at (0,.65) {2. Delete 30; promote 40};
\node[utpnode,minimum size=7mm] (n0) at (0,0) {50};
\node[utpnodehot,minimum size=7mm] (n1) at (-1.4,-1) {40};
\node[utpnode,minimum size=7mm] (n2) at (1.4,-1) {70};
\node[utpnode,minimum size=7mm] (n5) at (0.7,-2) {60};
\node[utpnode,minimum size=7mm] (n6) at (2.1,-2) {80};
\draw[utpedge] (n0)--(n1);
\draw[utpedge] (n0)--(n2);
\draw[utpedge] (n2)--(n5);
\draw[utpedge] (n2)--(n6);
\end{scope}\begin{scope}[xshift=5.8cm,yshift=-3.8cm]
\node[utplabel,anchor=south,color=UTPNavy] at (0,.65) {3. Delete 50; copy successor 60};
\node[utpnodehot,minimum size=7mm] (n0) at (0,0) {60};
\node[utpnode,minimum size=7mm] (n1) at (-1.4,-1) {40};
\node[utpnode,minimum size=7mm] (n2) at (1.4,-1) {70};
\node[utpnode,minimum size=7mm] (n6) at (2.1,-2) {80};
\draw[utpedge] (n0)--(n1);
\draw[utpedge] (n0)--(n2);
\draw[utpedge] (n2)--(n6);
\end{scope}
\draw[utparrow] (2.5,-.8)--(3.3,-.8);
\draw[utparrow] (2.5,-4.6)--(3.3,-4.6);
\node[utplabel,align=center] at (2.9,-6.6) {Final inorder: 40,60,70,80. Copying 60 also requires deleting its old node under 70.};
\end{tikzpicture}
\end{labgraphic}
\end{utpcompactlabfigure}


\lstinputlisting[firstline=31,lastline=65]{all_experiments/Lab08/bst.cpp}


\tcbbreak

\begin{utpsteps}\setcounter{enumi}{2}\item \textbf{Implement inorder, height, isBST and helpers.}\end{utpsteps}

\texttt{bst.cpp} --- lines 66--97.

\lstinputlisting[firstline=66,lastline=97]{all_experiments/Lab08/bst.cpp}


\tcbbreak

\begin{utpsteps}\setcounter{enumi}{3}\item \textbf{Implement draw, destroy, join and helpers.}\end{utpsteps}

\texttt{bst.cpp} --- lines 98--131.

\lstinputlisting[firstline=98,lastline=131]{all_experiments/Lab08/bst.cpp}


\tcbbreak

\begin{utpsteps}\setcounter{enumi}{4}\item \textbf{Connect the test cases in main.}\end{utpsteps}

\texttt{bst.cpp} --- lines 132--164.

% Worked operation using existing chapter styles.
\begin{utpcompactlabfigure}[At a matching node: return the subtree that replaces it]
\begin{labgraphic}
\begin{tikzpicture}[font=\small]

\node[utpstep,align=center] (q) at (5,0) {\texttt{value == root->data}};
\node[utpstep,align=left,text width=4.8cm] (a) at (1,-2) {\textbf{Zero or one child}\\Choose the existing child (or null).\\Delete root; return child.};
\node[utpstep,align=left,text width=5.3cm] (b) at (9,-2) {\textbf{Two children}\\Find leftmost node of right subtree.\\Copy its data; delete it below.\\Return the original root.};
\draw[utparrow] (q)--node[utplabel,above left,align=center]{left is null\\or right is null}(a);
\draw[utparrow] (q)--node[utplabel,above right]{both children exist}(b);
\node[utplabel,align=center] at (5,-3.5) {The caller assigns the returned pointer to its child link (or to the tree root).};
\end{tikzpicture}
\end{labgraphic}
\end{utpcompactlabfigure}


\textbf{Trace checkpoint for the supplied main.} Follow the three deletion cases in order: remove leaf 20; replace 30 by its only child 40; replace two-child root 50 by successor 60 and remove the successor from its former position. Assign the returned root after each call. Inorder must remain sorted, root must hold 60, and four live nodes must remain before destroy.

\lstinputlisting[firstline=132,lastline=164]{all_experiments/Lab08/bst.cpp}


\tcbbreak

\begin{utpsteps}\setcounter{enumi}{5}\item \textbf{Connect the test cases in main.}\end{utpsteps}

\texttt{bst.cpp} --- lines 165--181. Continue in the same file; do not add another header or opening brace.

\lstinputlisting[firstline=165,lastline=181]{all_experiments/Lab08/bst.cpp}


\tcbbreak

\textbf{Build, run, and inspect bst.}

From this program\textquotesingle s folder, run these commands in PowerShell:

\begin{lstlisting}[style=utpterminal,language={}]

g++ -std=c++17 -Wall -Wextra bst.cpp -o bst.exe

.\bst.exe

\end{lstlisting}

\textbf{Expected output excerpt.} Compare values, order, and pass/fail messages; spacing need not match.

\begin{lstlisting}[style=utpterminal,language={}]

== Insert 50, 30, 70, 20, 40, 60, 80 ==
            80
        70
            60
    50
            40
        30
            20
after insertion
  inorder: 20 30 40 50 60 70 80   height: 2
  [ok]   inorder output is sorted
  [ok]   every node respects the BST rule

== Search ==
  contains(60): yes   contains(65): no

\end{lstlisting}

The complete captured run is in \texttt{labs/outputs/bst.txt}.

\begin{tryit}

Insert a duplicate key and verify that the inorder output is unchanged. Delete the root from a one-node tree; assign the returned pointer and confirm it is null.

\end{tryit}


\tcbbreak
\textbf{Program 2: Max-heap and min-heap from start to finish.}

\textbf{File:} \texttt{heap.cpp}.

\textbf{Chapter connection:} Array representation of a heap; Insertion into a max-heap.

Insertion sifts up; root removal replaces the root with the last value and sifts down. Check child bounds before reading them. The min-heap uses the opposite comparisons, and both implementations are verified after mutations.

\begin{utpcompactlabfigure}[How this program runs]
\begin{labgraphic}
\begin{tikzpicture}[font=\small,execute at begin node={\hyphenpenalty=10000\exhyphenpenalty=10000}]
\node[utpstep,text width=2.50cm,minimum height=1.15cm] (s0) at (0.0,0) {\textbf{1}\\Append a key};
\node[utpstep,text width=2.50cm,minimum height=1.15cm] (s1) at (2.95,0) {\textbf{2}\\Sift upward};
\node[utpstep,text width=2.50cm,minimum height=1.15cm] (s2) at (5.9,0) {\textbf{3}\\Remove / sift downward};
\node[utpstep,text width=2.50cm,minimum height=1.15cm] (s3) at (8.85,0) {\textbf{4}\\Verify max and min};
\draw[utparrow] (s0.east)--(s1.west);
\draw[utparrow] (s1.east)--(s2.west);
\draw[utparrow] (s2.east)--(s3.west);
\end{tikzpicture}
\end{labgraphic}
\end{utpcompactlabfigure}


\begin{utpsteps}\setcounter{enumi}{0}\item \textbf{Headers and initial definitions.}\end{utpsteps}

\texttt{heap.cpp} --- lines 14--30.

\lstinputlisting[firstline=14,lastline=30]{all_experiments/Lab08/heap.cpp}


\tcbbreak

\begin{utpsteps}\setcounter{enumi}{1}\item \textbf{Implement popMaxHeap, pushMinHeap.}\end{utpsteps}

\texttt{heap.cpp} --- lines 31--61.

% Worked lab input; chapter fonts, colours and TikZ styles retained.
\begin{utpcompactlabfigure}[Sift-up and sift-down on the lab heap]
\begin{labgraphic}
\begin{tikzpicture}[font=\small]
\node[utpcell,minimum width=0.88cm,minimum height=.65cm] at (0.00,0.00) {50};
\node[utplabel] at (0.00,0.60) {0};
\node[utpcell,minimum width=0.88cm,minimum height=.65cm] at (0.92,0.00) {20};
\node[utplabel] at (0.92,0.60) {1};
\node[utpcell,minimum width=0.88cm,minimum height=.65cm] at (1.84,0.00) {30};
\node[utplabel] at (1.84,0.60) {2};
\node[utpcell,minimum width=0.88cm,minimum height=.65cm] at (2.76,0.00) {10};
\node[utplabel] at (2.76,0.60) {3};
\node[utphot,minimum width=0.88cm,minimum height=.65cm] at (3.68,0.00) {40};
\node[utplabel] at (3.68,0.60) {4};
\node[utplabel,anchor=west,align=left,text width=5.2cm] at (4.8,0.00) {Append 40: i=4, parent p=1};
\draw[utparrow] (-.85,-0.35)--(-.85,-0.85);
\node[utpcell,minimum width=0.88cm,minimum height=.65cm] at (0.00,-1.20) {50};
\node[utphot,minimum width=0.88cm,minimum height=.65cm] at (0.92,-1.20) {40};
\node[utpcell,minimum width=0.88cm,minimum height=.65cm] at (1.84,-1.20) {30};
\node[utpcell,minimum width=0.88cm,minimum height=.65cm] at (2.76,-1.20) {10};
\node[utpcell,minimum width=0.88cm,minimum height=.65cm] at (3.68,-1.20) {20};
\node[utplabel,anchor=west,align=left,text width=5.2cm] at (4.8,-1.20) {40 $>$ 20: swap 4 and 1\\50 $\ge$ 40: stop sifting up};
\draw[utparrow] (-.85,-1.55)--(-.85,-2.05);
\node[utphot,minimum width=0.88cm,minimum height=.65cm] at (0.00,-2.40) {20};
\node[utpcell,minimum width=0.88cm,minimum height=.65cm] at (0.92,-2.40) {40};
\node[utpcell,minimum width=0.88cm,minimum height=.65cm] at (1.84,-2.40) {30};
\node[utpcell,minimum width=0.88cm,minimum height=.65cm] at (2.76,-2.40) {10};
\node[utpdim,minimum width=0.88cm,minimum height=.65cm] at (3.68,-2.40) {--};
\node[utplabel,anchor=west,align=left,text width=5.2cm] at (4.8,-2.40) {Remove 50; move last value 20 to root\\Shorten active size to 4};
\draw[utparrow] (-.85,-2.75)--(-.85,-3.25);
\node[utphot,minimum width=0.88cm,minimum height=.65cm] at (0.00,-3.60) {40};
\node[utphot,minimum width=0.88cm,minimum height=.65cm] at (0.92,-3.60) {20};
\node[utpcell,minimum width=0.88cm,minimum height=.65cm] at (1.84,-3.60) {30};
\node[utpcell,minimum width=0.88cm,minimum height=.65cm] at (2.76,-3.60) {10};
\node[utpdim,minimum width=0.88cm,minimum height=.65cm] at (3.68,-3.60) {--};
\node[utplabel,anchor=west,align=left,text width=5.2cm] at (4.8,-3.60) {Choose larger child 40; swap with 20\\20 $\ge$ 10: stop sifting down};
\end{tikzpicture}
\end{labgraphic}
\end{utpcompactlabfigure}


\lstinputlisting[firstline=31,lastline=61]{all_experiments/Lab08/heap.cpp}


\tcbbreak

\begin{utpsteps}\setcounter{enumi}{2}\item \textbf{Implement popMinHeap, isMaxHeap, isMinHeap and helpers.}\end{utpsteps}

\texttt{heap.cpp} --- lines 62--98.

\lstinputlisting[firstline=62,lastline=98]{all_experiments/Lab08/heap.cpp}


\tcbbreak

\begin{utpsteps}\setcounter{enumi}{3}\item \textbf{Connect the test cases in main.}\end{utpsteps}

\texttt{heap.cpp} --- lines 99--135.

\lstinputlisting[firstline=99,lastline=135]{all_experiments/Lab08/heap.cpp}


\tcbbreak

\begin{utpsteps}\setcounter{enumi}{4}\item \textbf{Connect the test cases in main.}\end{utpsteps}

\texttt{heap.cpp} --- lines 136--148. Continue in the same file; do not add another header or opening brace.

\lstinputlisting[firstline=136,lastline=148]{all_experiments/Lab08/heap.cpp}


\tcbbreak

\textbf{Build, run, and inspect heap.}

From this program\textquotesingle s folder, run these commands in PowerShell:

\begin{lstlisting}[style=utpterminal,language={}]

g++ -std=c++17 -Wall -Wextra heap.cpp -o heap.exe

.\heap.exe

\end{lstlisting}

\textbf{Expected output excerpt.} Compare values, order, and pass/fail messages; spacing need not match.

\begin{lstlisting}[style=utpterminal,language={}]

== Max-heap: insert 30, 10, 50, 20, 40 ==
  insert 30 -> [30]
  insert 10 -> [30 10]
  insert 50 -> [50 10 30]
  insert 20 -> [50 20 30 10]
  insert 40 -> [50 40 30 10 20]
  [ok]   after every insertion no parent is smaller than a child
  [ok]   final max-heap is [50 40 30 10 20]

== Max-heap: delete the root twice ==
  removed 50 -> [40 20 30 10]
  [ok]   heap property restored after deletion
  removed 40 -> [30 20 10]
  [ok]   heap property restored after deletion

\end{lstlisting}

The complete captured run is in \texttt{labs/outputs/heap.txt}.

\begin{tryit}

Drain both heaps completely. Max-heap removals must descend; min-heap removals must ascend. After each removal, check every parent-child pair, then test one more pop on empty.

\end{tryit}

\textbf{Instructor-led B-tree extension.} Use the degree and occupancy rules specified in class to draw split, borrow, and merge. The complete programs here cover BSTs, height checks, and min/max heaps; B-tree implementation is outside these supplied program files.
\end{labproject}

\begin{verification}
$\square$ Every listed program compiles independently with its own \texttt{main()}.\\
$\square$ Recorded outputs match the traces; no built-in check prints \texttt{[FAIL]}.\\
$\square$ Explain one code-to-diagram state change and demonstrate one additional test.
\end{verification}
